\documentclass[english,hebrew]{article}
\usepackage[HE8,T1]{fontenc}
\usepackage[utf8]{inputenc}
\setcounter{secnumdepth}{-2}
\setlength{\parindent}{0bp}
\usepackage{amsmath}
\usepackage{amsthm}
\usepackage{amssymb}

\makeatletter
\theoremstyle{definition}
\newtheorem*{problem*}{\protect\problemname}
\theoremstyle{definition}
\newtheorem*{sol*}{\protect\solutionname}

\@ifundefined{date}{}{\date{}}
\usepackage{graphicx}
\usepackage{geometry}
\newcommand{\limi}{\underset{n\rightarrow\infty}{\lim}}
\newcommand{\limxz}{\underset{x\rightarrow x_{0}}{\lim}}
\newcommand{\limfxz}{\underset{x\rightarrow x_{0}}{\lim}f(x)}
\newcommand{\limfxm}{\underset{x\rightarrow x_{0}^{-}}{\lim}f(x)}
\newcommand{\limfxp}{\underset{x\rightarrow x_{0}^{+}}{\lim}f(x)}
\newcommand{\limxm}{\underset{x\rightarrow x_{0}^{-}}{\lim}}
\newcommand{\limxp}{\underset{x\rightarrow x_{0}^{+}}{\lim}}
\newcommand{\sqrm}{M_{m\times m}(\mathbb{F})}
\newcommand{\seqa}{(a_{n})_{n=1}^{\infty}}
\newcommand{\seqx}{(x_{n})_{n=1}^{\infty}}
\newcommand{\funcdr}{f:D\rightarrow\mathbb{R}}
\newcommand{\der}{\limxz\frac{f(x)-f(x_{0})}{x-x_{0}}}
\usepackage[all]{pdfprivacy}

\makeatother

\usepackage{babel}
\addto\captionsenglish{\renewcommand{\problemname}{Problem}}
\addto\captionsenglish{\renewcommand{\solutionname}{Solution}}
\addto\captionshebrew{\renewcommand{\problemname}{בעיה}}
\addto\captionshebrew{\renewcommand{\solutionname}{פתרון}}
\providecommand{\problemname}{בעיה}
\providecommand{\solutionname}{פתרון}

\begin{document}
\newgeometry{top=1cm, bottom=2cm, left=1cm, right=1cm}
\title{סמסטר {\beginL 2027\endL}א, בוחן {\beginL 1\endL}}
\date{{\beginL 30.10.2026\endL}}
\maketitle
\begin{abstract}
בוחן {\beginL 1\endL} בקורס אלגברה לינארית 2, סמסטר תשפ”ז א’. השאלות
נכתבו על ידי צוות הקורס באוניברסיטה העברית. הפתרון הוא שלי ולא רשמי.
\end{abstract}
\begin{problem*}
נתונה מטריצה \L{$A$}: \L{%
\[
\begin{pmatrix}-1 & 0 & 2\\
0 & 2 & 0\\
-4 & 0 & 5
\end{pmatrix}
\]
}%

א. מצאו את הערכים עצמיים של \L{$A$}.

ב. מצאו את הוקטורים העצמיים המתאימים לכל ערך עצמי.

ג. הוכיחו כי מטריצה \L{$A$} לכסינה, מצאו מטריצה הפיכה \L{$Q$} ומטריצה
אלכסונית \L{$D$} כך ש-\L{$A=QDQ^{-1}$}.

ד. הגדירו מושג: \textquotedblright תת-מרחב שמור של אופרטור לינארי\textquotedblright .
\end{problem*}
\begin{sol*}
א. נבדוק מתי הפולינום האופייני של המטריצה מתאפס.\L{%
\[
\det\left(\lambda I-A\right)=0
\]
}%
\L{%
\[
\iff\det\left(\begin{array}{ccc}
\lambda+1 & 0 & -2\\
0 & \lambda-2 & 0\\
4 & 0 & \lambda-5
\end{array}\right)=0
\]
}%
\L{%
\[
\iff\begin{vmatrix}\lambda+1 & 0 & -2\\
0 & \lambda-2 & 0\\
4 & 0 & \lambda-5
\end{vmatrix}=0
\]
}%
\L{%
\[
\iff\left(\lambda+1\right)\cdot\begin{vmatrix}\lambda-2 & 0\\
0 & \lambda-5
\end{vmatrix}+\left(-2\right)\cdot\begin{vmatrix}0 & \lambda-2\\
4 & 0
\end{vmatrix}=0
\]
}%
\L{%
\[
\iff\left(\lambda+1\right)\left(\lambda-2\right)\left(\lambda-5\right)+8\left(\lambda-2\right)=0
\]
}%
\L{%
\[
\iff\left(\lambda-2\right)\left(\left(\lambda+1\right)\left(\lambda-5\right)+8\right)=0
\]
}%
\L{%
\[
\iff\left(\lambda-2\right)\left(\lambda^{2}-4\lambda+3\right)=0
\]
}%
\L{%
\[
\iff\left(\lambda-2\right)\left(\lambda-3\right)\left(\lambda-1\right)=0
\]
}%

ולפי המשפט על אפיון ערכים עצמיים באמצעות הפולינום האופייני נקבל שהערכים
העצמיים הם \L{$2,3,1$}.

ב. אנו יודעים שעבור וקטור עצמי \L{$v_{\lambda}$} של ערך עצמי \L{$\lambda$}
מתקיים \L{$A\cdot v_{\lambda}=\lambda I\cdot v_{\lambda}$}, כלומר
\L{$\left(\lambda I-A\right)v_{\lambda}=0$}.

נציב כל אחד מהערכים העצמיים, נציג וקטור כללי ונקבל בסיס לתמ\textquotedblright ו
שמשרה כל ערך עצמי.

\L{$\lambda=1$}:\L{%
\[
\begin{pmatrix}2 & 0 & -2\\
0 & -1 & 0\\
4 & 0 & -4
\end{pmatrix}\begin{pmatrix}x\\
y\\
z
\end{pmatrix}=\begin{pmatrix}0\\
0\\
0
\end{pmatrix}
\]
}%
\L{%
\[
\iff\begin{cases}
2x-2z & =0\\
-y & =0\\
4x-4z & =0
\end{cases}
\]
}%
\L{%
\[
\iff x=z,\quad y=0
\]
}%

נציב \L{$x=1$} ונקבל\L{%
\[
E_{\lambda=1}=\text{Span}\begin{pmatrix}1\\
0\\
1
\end{pmatrix}
\]
}%

כאשר הוקטורים העצמיים של \L{$\lambda=1$} הם \L{$v\in E_{\lambda=1}\setminus\left\{ 0_{V}\right\} $}.

\L{$\lambda=3$}:\L{%
\[
\begin{pmatrix}4 & 0 & -2\\
0 & 1 & 0\\
4 & 0 & -2
\end{pmatrix}\begin{pmatrix}x\\
y\\
z
\end{pmatrix}=\begin{pmatrix}0\\
0\\
0
\end{pmatrix}
\]
}%
\L{%
\[
\iff\begin{cases}
4x-2z & =0\\
y & =0\\
4x-2z & =0
\end{cases}
\]
}%
\L{%
\[
\iff2x=z,\quad y=0
\]
}%

נציב \L{$x=1$} ונקבל\L{%
\[
E_{\lambda=3}=\text{Span}\begin{pmatrix}1\\
0\\
2
\end{pmatrix}
\]
}%

כאשר הוקטורים העצמיים של \L{$\lambda=3$} הם \L{$v\in E_{\lambda=3}\setminus\left\{ 0_{V}\right\} $}.

\L{$\lambda=2$}:\L{%
\[
\begin{pmatrix}3 & 0 & -2\\
0 & 0 & 0\\
4 & 0 & -3
\end{pmatrix}\begin{pmatrix}x\\
y\\
z
\end{pmatrix}=\begin{pmatrix}0\\
0\\
0
\end{pmatrix}
\]
}%
\L{%
\[
\iff\begin{cases}
3x-2z & =0\\
4x-3z & =0
\end{cases}
\]
}%
\L{%
\[
\iff x=z=0
\]
}%

נציב \L{$y=1$} ונקבל\L{%
\[
E_{\lambda=2}=\text{Span}\begin{pmatrix}0\\
1\\
0
\end{pmatrix}
\]
}%

כאשר הוקטורים העצמיים של \L{$\lambda=2$} הם \L{$v\in E_{\lambda=2}\setminus\left\{ 0_{V}\right\} $}.

ג. הראינו שיש ל-\L{$A$} שלושה ערכים עצמיים שונים ומתקיים \L{$\dim\left(\mathbb{F}^{3}\right)=3$},
לכן לפי המשפט שאומר שאם לאופרטור על מ\textquotedblright ו ממימד \L{$n$}
יש \L{$n$} ערכים עצמיים שונים, האופרטור לכסין, נקבל שהמטריצה \L{$A$}
היא לכסינה. 

לפי אותו משפט והאלגוריתם שלמדנו, \L{$A$} היא לכסינה כאשר המטריצה
האלכסונית הדומה לה היא המטריצה עם הערכים העצמיים על האלכסון:\L{%
\[
D=\begin{pmatrix}1 & 0 & 0\\
0 & 3 & 0\\
0 & 0 & 2
\end{pmatrix}
\]
}%
 ומטריצת מעבר בסיס היא מטריצה שעמודותיה הן וקטורים שפורשים את התמ\textquotedblright וים
שמשרים הערכים העצמיים בהתאמה\L{%
\[
Q=\begin{pmatrix}1 & 1 & 0\\
0 & 0 & 1\\
1 & 2 & 0
\end{pmatrix}
\]
}%

ד. יהיו \L{$V$} מ\textquotedblright ו, \L{$T\in\mathcal{L}\left(V\right)$}.
אז \L{$U\subseteq V$} תמ\textquotedblright ו שמור תחת \L{$T$} אם\L{%
\[
\forall u\in U:T\left(u\right)\in U
\]
}%
\end{sol*}

\end{document}
